\section{Despliegue, CI y observabilidad}
\label{sec:despliegue}
\textit{Responsables: Fabricio (despliegue y alarmas, PLT-01, PLT-03, PLT-05, PLT-06) y Mariano (CI,
PLT-02; dashboard, PERF-01).}

\subsection{Infraestructura como código}
Todo se declara con Serverless Framework v4, con un \texttt{serverless.yml} por stack y los stages
\texttt{dev}, \texttt{staging} y \texttt{prod}. Como el Learner Lab no deja crear roles, todas las
funciones usan \texttt{arn:aws:iam::\$\{aws:accountId\}:role/LabRole}. Los stacks se conectan por
parámetros SSM (\texttt{/nubeduca/<stage>/...}) y \textit{exports} de CloudFormation; nadie copia
ARNs a mano.

\begin{table}[htbp]
\centering\small
\begin{tabularx}{\textwidth}{lXl}
\toprule
Stack & Qué crea & Repositorio \\
\midrule
\texttt{infra/base} & Bucket de despliegue y parámetros SSM compartidos & \texttt{nubeduca-backend} \\
\texttt{infra/bus} & Bus \texttt{lms-<stage>}, Archive y SNS FIFO de notas & \texttt{nubeduca-backend} \\
\texttt{infra/auth} & User Pool y App Client de Cognito & \texttt{nubeduca-backend} \\
\texttt{infra/frontend} & App de Amplify y variables \texttt{VITE\_*} desde SSM & \texttt{nubeduca-backend} \\
\texttt{infra/despliegue} & VM \texttt{nubeduca-despliegue} (Session Manager) & \texttt{nubeduca-backend} \\
\texttt{infra/observabilidad} & Dashboard \texttt{nubeduca-<stage>} de CloudWatch & \texttt{nubeduca-backend} \\
Cada microservicio & API Gateway, Lambdas, tabla, Pipe, regla, cola, DLQ y alarma & \texttt{nubeduca-ms-*} \\
\bottomrule
\end{tabularx}
\caption{Stacks de Serverless Framework.}
\label{tab:stacks}
\end{table}

\subsection{Despliegue completo con \texttt{desplegar.sh}}
\texttt{scripts/desplegar.sh <stage>} despliega todo en la cuenta del \texttt{.env}, en orden:
\texttt{infra/base} $\rightarrow$ \texttt{infra/bus} $\rightarrow$ \texttt{infra/auth}
$\rightarrow$ cada \texttt{nubeduca-ms-*} (clonado desde \texttt{main}) $\rightarrow$ frontend.
Antes de crear nada, \texttt{scripts/cuenta.sh} verifica que las credenciales sean de la cuenta
esperada. Al final hace una prueba de humo (\texttt{GET /health} de cada API) y un resumen. Opciones:
\texttt{--con-seed} (carga los 30\,000 registros por servicio), \texttt{--solo <ms>} y
\texttt{--sin-frontend}.

\pendiente{Fabricio: describir la VM \texttt{nubeduca-despliegue} (rol de instancia
\texttt{LabInstanceProfile}, sin puertos abiertos) y la cuenta de integración oficial (D2) cuando se
defina.}

\subsection{Integración continua}
Cada repositorio corre GitHub Actions en cada PR: lint (ruff, yamllint) y pruebas con moto (sin AWS).
\texttt{nubeduca-contratos} valida además los esquemas y su compatibilidad con \texttt{main}. Ramas
\texttt{<tipo>/<ID>-descripcion}, PR a \texttt{main} y \textit{squash merge}.

\pendiente{Indicar si se activó la protección de \texttt{main} con checks obligatorios (requiere un
admin de la organización). Mientras no esté activa, el check en rojo avisa pero no impide el merge.}

\subsection{Observabilidad}
\begin{itemize}
  \item \textbf{CloudWatch Logs}: cada Lambda escribe logs estructurados
        (\texttt{\{"tipo": "INFO"|"ERROR", "log\_datos": \{...\}\}}).
  \item \textbf{Alarmas}: una alarma por DLQ (dispara con 1 mensaje). \texttt{scripts/verificar\_alarmas.sh
        <stage>} lista su estado.
  \item \textbf{X-Ray}: \textit{tracing} activo en API Gateway y Lambda. El visor no se puede usar en el
        Learner Lab (sección~\ref{sec:limitaciones}); la evidencia de la traza son los \texttt{TraceId}
        en CloudWatch Logs.
  \item \textbf{Dashboard} \texttt{nubeduca-<stage>} (PERF-01): latencia p95 y errores 4XX/5XX de API
        Gateway, duración, errores, throttles y concurrencia de Lambda, throttles de DynamoDB, mensajes y
        antigüedad en SQS, DLQ y reglas de EventBridge. Usa expresiones \texttt{SEARCH}, así que cada
        microservicio nuevo aparece sin editar el dashboard.
\end{itemize}

\subsection{Resultados y evidencias}
\paragraph{Staging en la cuenta de Dev 1.}
\pendiente{Fabricio: resultado de \texttt{desplegar.sh staging --sin-frontend} en su cuenta y, cuando
se defina, en la cuenta de integración.}

\paragraph{Staging en la cuenta de Dev 2 (21-Nov-2026).}
Con todos los PRs del ADR-008 ya mergeados, \texttt{scripts/desplegar.sh staging --sin-frontend}
desplegó \texttt{infra/base}, \texttt{infra/bus}, \texttt{infra/auth} y los seis microservicios, sin
pendientes ni fallos:
\begin{itemize}
  \item \texttt{GET /health} respondió 200 en los seis microservicios.
  \item \texttt{verificar\_alarmas.sh staging}: las 9 alarmas de DLQ en \texttt{OK}.
  \item Los cinco microservicios que publican eventos tienen su Pipe
        \texttt{nubeduca-ms-*-staging-outbox} en \texttt{RUNNING} (ms-notificaciones no publica).
\end{itemize}
Al terminar se borró staging de esa cuenta para no consumir presupuesto.

\begin{lstlisting}
== Prueba de humo: GET /health
  ms-usuarios: 200        ms-cursos: 200          ms-contenidos: 200
  ms-evaluaciones: 200    ms-calificaciones: 200  ms-notificaciones: 200

== Resumen (staging)
  Desplegados: infra/base infra/bus infra/auth ms-usuarios ms-cursos
               ms-contenidos ms-evaluaciones ms-calificaciones ms-notificaciones
  Pendientes:  ninguno
  Fallidos:    ninguno

$ scripts/verificar_alarmas.sh staging
nubeduca-staging-calificaciones-auditoria-dlq    OK
nubeduca-staging-calificaciones-columnas-dlq     OK
nubeduca-staging-calificaciones-eventos-dlq      OK
nubeduca-staging-calificaciones-notas-dlq        OK
nubeduca-staging-contenidos-eventos-dlq          OK
nubeduca-staging-cursos-eventos-dlq              OK
nubeduca-staging-evaluaciones-eventos-dlq        OK
nubeduca-staging-notificaciones-eventos-dlq      OK
nubeduca-staging-notificaciones-notas-dlq        OK
Todas en verde (OK) en staging.

$ aws pipes list-pipes   (Pipes de staging)
nubeduca-ms-calificaciones-staging-outbox   RUNNING
nubeduca-ms-contenidos-staging-outbox       RUNNING
nubeduca-ms-cursos-staging-outbox           RUNNING
nubeduca-ms-evaluaciones-staging-outbox     RUNNING
nubeduca-ms-usuarios-staging-outbox         RUNNING
\end{lstlisting}
\evidencia{dashboard-api.png}{Dashboard \texttt{nubeduca-dev} durante PERF-01: latencia p95 de API
Gateway (el listado de cursos ronda 1{,}8\,s) y sin errores 4XX/5XX.}{fig:dashboard-api}
\evidencia{dashboard-lambda-dynamo.png}{Dashboard \texttt{nubeduca-dev} durante PERF-01: concurrencia
de Lambda y lecturas consumidas de DynamoDB, altas porque el listado recorre toda la partición del
tenant.}{fig:dashboard-dynamo}
